\chapter{Sticker systems and molecular memory}
A sequence need not change for its stored information to change. Imagine
four address regions on one long DNA strand. Short complementary strands
occupy addresses 1 and 3; addresses 0 and 2 remain exposed. Under an
occupancy convention, that complex represents $0101$ when bits are printed
in increasing address order. Removing the sticker at address 1 changes the
register without shortening the backbone.

Can this representation compute $c=a\land b$ while preserving $a$ and $b$?
Can it reuse $c$? If two distinct registers become identical after a clear,
have molecules disappeared, or only distinguishable information?
Those questions lead from representation to a small executable machine.

Chapter 13 used fixed assignment sequences and selection. This chapter
keeps a common backbone and changes occupancy. We assume familiarity with
antiparallel pairing, strand domains, multiset counts and ideal separation.
Our exact model describes logical operations; it does not prescribe
selective sticker removal in a laboratory.

\section{Separate the memory strand from its current contents}
\index{occupancy register}
The sticker model distinguishes a long \Term{memory strand}, address-specific
short \Term{stickers}, and a \Term{memory complex} consisting of the strand
plus its bound stickers \cite{stickers}. We use $w$ non-overlapping domains
$D_0,\ldots,D_{w-1}$ as addresses. Bit $b_i=1$ means its complementary
sticker is present; $b_i=0$ means it is absent. One address is a designed
sequence region, not a single nucleotide.

\Plate{01-memory}{A four-address memory complex stores occupancy 0101
in increasing address order. The long backbone is continuous and runs
$5'$ to $3'$; each separate sticker runs antiparallel. Dashed rungs denote
pairing, not covalent joins.}{fig:memory}

The physical and computational objects must not be conflated. In the
software reference, the state is an integer
\begin{equation}
 s=\sum_{i=0}^{w-1}b_i2^i,\qquad 0\leq s<2^w.
 \label{eq:register}
\end{equation}
Address 0 is the least significant bit. The printed order $b_0b_1b_2$
is therefore the reverse of conventional most-significant-bit-first
binary formatting. We write explicit tuples whenever that distinction
could confuse a trace.

A tube is a multiset $P(s)$ of memory complexes. Two complexes can have
identical backbones and different occupancies. Two identical occupancies
can have many physical copies. Sequence identity, logical state and copy
identity are three different notions.

\section{A write is a map on states and a pushforward on counts}
\index{selective clear}\index{copy conservation}
At one address $i$, the ideal set and clear maps are
\begin{equation}
 S_i(s)=s\mathbin{\mathrm{OR}}2^i,\qquad
 C_i(s)=s\mathbin{\mathrm{AND}}\neg 2^i.
 \label{eq:write}
\end{equation}
The second expression is evaluated on a fixed-width register; equivalently,
subtract $2^i$ only when bit $i$ is one. Python's unbounded signed complement
is safe here because the validated nonnegative input has no higher set bits.

For any deterministic map $f$, the output population is
\begin{equation}
 (f_*P)(t)=\sum_{s:f(s)=t}P(s).
 \label{eq:pushforward}
\end{equation}
This is a \Term{pushforward}: all input states mapping to the same output
contribute their counts. A dictionary assignment such as
$Q[f(s)]\leftarrow P(s)$ is wrong when states collide. It overwrites one
count instead of adding it.

\Listing{write}
\Plate{02-write}{Set associates a complementary sticker with a vacant
address; clear removes that address's sticker in the formal model.
The scaffold and other addresses remain unchanged. The clear arrow states
an assumed operation, not a validated address-selective melting recipe.}{fig:write}

Logical identities help test the implementation:
\begin{align}
 S_iS_i&=S_i,& C_iC_i&=C_i,\\
 C_iS_i&=C_i,& S_iC_i&=S_i.
 \label{eq:algebra}
\end{align}
Composition acts rightmost first. Different addresses commute. At the
same address, set and clear generally do not commute: their final states
are determined by the last write. Idempotence means repeating a logical
write changes no additional bits; it does not mean adding reagent twice
has no chemical or economic consequences.

Equation~\ref{eq:pushforward} conserves the number of memory backbones:
$\sum_t(f_*P)(t)=\sum_sP(s)$. It does not conserve the number of bound
stickers. Setting empty addresses consumes free stickers; clearing releases
or otherwise removes them from the bound population. A chemical mass
balance must include those pools and any removal mechanism.

\section{Separation turns a uniform write into conditional logic}
\index{AND gate}\index{address aliasing}
A uniform set affects every complex in its tube. Conditional computation
therefore separates the population, writes only one branch, and combines
the branches. Ideal separation at address $i$ produces
\begin{equation}
 P_1(s)=P(s)\mathbf1[b_i=1],\qquad
 P_0(s)=P(s)\mathbf1[b_i=0].
\end{equation}
Combining tubes adds multiplicities. It is not set union with duplicates
discarded. As in Chapter 13, separation partitions one available inventory;
it is not duplication.

\Listing{separate}

To compute $c:=a\land b$ at three distinct addresses:
first clear $c$ across the population. Separate on $a$.
Separate the $a=1$ branch on $b$.
Set $c$ only in the $a=b=1$ branch, then merge the three disjoint branches.
The input bits are never written.

\Plate{03-and}{Conditional writing implements an AND gate over an entire
population. Each row begins with five copies of the labeled input register.
Only the 11 branch acquires the output sticker; all twenty memory
backbones return to the final tube.}{fig:and}

\Listing{and_gate}
\begin{center}\input{\Generated/gate-table.tex}\end{center}

Why clear the output first? Without that step, the algorithm computes
$c_{\mathrm{new}}=c_{\mathrm{old}}\lor(a\land b)$, not assignment to AND.
A stale output bit is a state-initialization bug, not an error in the truth
table. The reference tests every initial value of $c$, every input pair
and every permutation of the three address roles.

Why forbid $c=a$ or $c=b$? Clearing the output would destroy an input
before it was tested. An in-place gate needs a different schedule or
scratch storage. Address aliasing changes the algorithm even when every
individual operation remains well-defined.

\section{Reset destroys distinctions, not backbones}
\index{reset!information loss}
Consider three copies of occupancy $(0,0,0)$ and five of $(0,0,1)$.
Clearing address 2 maps both to $(0,0,0)$, yielding eight copies of one
logical state. The species count falls from two to one; the copy count
remains eight.

\Plate{04-reset}{Clear is many-to-one. Without a retained history tag,
the output does not identify which initial occupancy produced a complex.
A history address can preserve the distinction, at the cost of additional
storage and operations.}{fig:reset}

For a random initial bit $B$ with $\Pr(B=1)=p$, the uncertainty is
\begin{equation}
 H(B)=-p\log_2p-(1-p)\log_2(1-p)
\end{equation}
bits, taking $0\log_20=0$. A deterministic clear produces a constant bit
with entropy zero. For $p=1/2$, one bit of distinguishability is lost.
This information statement does not specify heat dissipated by the actual
DNA process: that would require a physical reset protocol, its environment
and thermodynamic assumptions. Logical irreversibility and measured
energetic cost are related questions, not interchangeable measurements.

A history tag makes the larger map injective: $(b,0)\mapsto(0,b)$ retains
the old bit elsewhere. Eventually clearing the history restores the
original problem. Scratch memory postpones disposal of information; it
does not make disposal disappear.

Reusable also does not mean indefinitely stable. A backbone may remain
intact while a sticker is lost, an address becomes unavailable, or an
off-target strand occupies it. A storage claim should specify retention
time, rewrite count, read disturbance, error definition and environment.

\section{Build an error budget before scaling the register}
\index{Markov error model}\index{memory!reuse}
Let $p$ be the probability that one required operation preserves a correct
memory complex, conditioned on correctness so far. With the same $p$
at each of $k$ operations and no recovery,
\begin{equation}
 \Pr(\text{still correct after }k)=p^k.
 \label{eq:retention}
\end{equation}
An independent equal-probability model implies this form; more generally,
a chain of conditional success probabilities gives $\prod_jp_j$.
The reference uses assigned values, not measurements.

\Listing{retained}
\Plate{05-reliability}{Computed survival under two illustrative constant
per-operation probabilities. The curve counts complexes that never leave
the correct trajectory; it is not a prediction of experimentally measured
sticker fidelity or of error repair.}{fig:reuse}

\begin{center}\input{\Generated/reuse-table.tex}\end{center}

For a richer model, separate false setting from false clearing. If
$\alpha$ is the probability that an empty address becomes incorrectly
occupied, and $\beta$ is the probability that an occupied address loses
its sticker during one storage interval, a row-vector distribution evolves as
\begin{equation}
 [p_0',p_1']=[p_0,p_1]
 \begin{bmatrix}1-\alpha&\alpha\\ \beta&1-\beta\end{bmatrix}.
 \label{eq:transition}
\end{equation}
The rows sum to one. When $\alpha+\beta>0$, a stationary occupancy
probability is $\pi_1=\alpha/(\alpha+\beta)$. This is a property of the
assumed Markov model, not necessarily an equilibrium binding fraction.
The latter needs concentrations, free energies and mass balance from
Chapter 6. In particular, repeated observations of the same address can
be correlated, so taking a majority vote is not automatically independent
error reduction.

\section{Modern molecular memory: compare the mechanisms}
The 1998 sticker paper separates abstract operations from proposed physical
implementations \cite{stickers}. Our reference adopts ideal addressed writes;
selective clear remains a physical obligation. It would be misleading to
infer arbitrary reliable rewrites from a correct integer bitmask program.

A 2026 Scaffolded DNA Computer study instead engineers a preferred
equilibrium configuration and reports renewable computations \cite{sdc}.
Its renewal experiments add input and blocker strands and reanneal.
This is not a free, address-independent implementation of our clear
primitive. The study also reports a substantial scale-dependent distinction:
short programs can use fast anneals, whereas demonstrated larger systems
require much longer anneals.

The deeper comparison is useful: a register interface specifies what a
programmer may assume; an energy landscape specifies how a molecular
population may realize a result. To connect them, measure cross-address
interference, time to a usable output, consumables, carryover and readout
accuracy across repeated inputs. A common use of DNA scaffolds does not
make two computational models equivalent.

\section{Exercises with worked answers}
\begin{enumerate}
\item \textbf{Representation.} Encode $(b_0,b_1,b_2)=(1,0,1)$.
\textit{Answer:} $s=1+0+4=5$. For $(0,1,1)$ the integer is 6;
the printed address-order tuple must not be read as conventional binary 011.
\item \textbf{Write.} Set address 1 in state 5, then clear address 0.
\textit{Answer:} $5\mathbin{\mathrm{OR}}2=7$; clearing the unit bit gives 6.
Only occupancy changes in the abstraction.
\item \textbf{Inventory.} Clear bit 2 from $\{0:3,4:5\}$.
\textit{Answer:} The output is $\{0:8\}$. Assigning rather than adding
dictionary counts would incorrectly retain only one contribution.
\item \textbf{Algebra.} Prove idempotence of setting.
\textit{Answer:} Once bit $i$ is one, OR with $2^i$ cannot change it.
All other bits are untouched, so $S_i(S_i(s))=S_i(s)$ for every register.
\item \textbf{Counterexample.} Show set and clear need not commute.
\textit{Answer:} Starting at zero, $S_0(C_0(0))=1$ but
$C_0(S_0(0))=0$. On distinct addresses, the two affected coordinates
are disjoint and the operations commute.
\item \textbf{Gate.} Omit the initial clear from AND with input $(0,1,1)$.
\textit{Answer:} The complex takes the $a=0$ branch and retains $c=1$.
Correct assignment requires $c=0$.
\item \textbf{Aliasing.} Explain why output address $c=a$ is forbidden.
\textit{Answer:} The initial clear forces $a=0$ before separation;
all complexes then bypass the only setting branch.
\item \textbf{Reliability.} At $p=0.99$, what is the 100-operation success
probability? \textit{Answer:} $0.99^{100}\approx0.366032$.
A one-percent per-step failure can be unacceptable in a long uncorrected run.
\item \textbf{Markov model.} For $\alpha=0.02$, $\beta=0.08$, compute
stationary occupancy. \textit{Answer:} $\pi_1=0.2$; check that
$[0.8,0.2]$ multiplied by Equation~\ref{eq:transition}'s matrix is unchanged.
This does not identify a molecular binding constant.
\item \textbf{Information.} Does clearing eight identical copies erase
three bits because $\log_2 8=3$? \textit{Answer:} Not from that fact alone.
Copy count is not the uncertainty in the logical state. Entropy requires
a probability distribution over possible values and a specified observer.
\item \textbf{Testing.} Extend the reference with XOR.
\textit{Rubric:} Preserve input addresses, clear the output, partition all
four cases without duplicated copies, set exactly 01 and 10, test arbitrary
initial output and multiplicities, and reject address aliasing.
\item \textbf{Experimental design.} Test selective clear.
\textit{Rubric:} Track intended removal, off-address loss, scaffold recovery,
released-strand carryover, rewrite accuracy and repeated-cycle readout.
Include no-clear controls and independently identifiable addresses.
\end{enumerate}

\section{What changes in the next model?}
\textbf{Reproduce the chapter.} From the repository root, run
\texttt{python3 drafts/ch14/build.py --assets-only}, then
\texttt{python3 -m pytest tests/test\_ch14\_reference.py}.
The first command regenerates \texttt{results.json}, tables, plots and literal
code listings; the tests compare the implementation with independent
oracles and explicit failure cases. Run the builder without
\texttt{--assets-only} to compile this review PDF.

We have implemented finite occupancy memory with exact copy accounting.
Set and clear change a register; separation makes a write conditional;
combination reunites inventories; resetting a bit can merge logical states.

\textbf{Working vocabulary.} The memory strand is the address-bearing
backbone; the memory complex includes its current stickers; pushforward
adds counts under a state transformation; aliasing means two roles share
an address; a reset is many-to-one unless history is retained.
Chapter 15 changes the representation again: splicing and
insertion--deletion systems transform the strings themselves. We will need
to track sequence boundaries and rewrite rules, not merely occupancy.
